% =============================================================================
% Ready-to-paste LaTeX: Experiment 4 -- reliability of the emotion ratings
% (Set 1 vs Set 2) and the ceiling it places on the audio -> emotion regression.
% Numbers from experiments/features/exp4_rating_reliability.py.
%
% Uses bib keys already in library.bib: eerola2011comparison, kang2025there.
% Placement: a subsection in Results just BEFORE the emotion-regression results
% (it sets the reference point against which R^2 = 0.56 should be read), or in
% Discussion as a limitations/validity subsection.
% =============================================================================

\subsection{Reliability of the Emotion Ratings}
\label{subsec:rating_reliability}

The emotion ratings are the target of the first pipeline stage, so their reliability
determines how well any audio model could possibly perform. The dataset of
\citet{eerola2011comparison} permits a direct estimate: its second set re-rates excerpts
drawn from the first with a different listener panel, giving two independent measurements
of the same musical material. The correspondence between the sets is carried by a
reference column in the second set rather than by clip number, which the two sets assign
independently; the mapping was verified acoustically by full-lag normalised
cross-correlation between each pair of audio files (median $0.964$, against $0.027$ for
randomly paired clips), and the single excerpt whose audio did not match its reference was
excluded. Seven excerpts occur twice within the second set, presented twice to the same
panel; averaging these leaves $102$ matched excerpts from $38$ soundtracks.

\begin{table}[t]
  \centering
  \caption{Agreement between the two listener panels on the same $102$ excerpts.
    ICC(C,1) treats a constant offset between panels as irrelevant (do they rank
    excerpts alike?), ICC(A,1) counts it as disagreement (do they give the same
    numbers?). Bias is the mean difference in rating points on the $1$--$9$ scale.}
  \label{tab:rating_reliability}
  \begin{tabular}{lcccc}
    \hline
    Emotion & $r$ & ICC(C,1) & ICC(A,1) & Bias \\
    \hline
    Valence & $0.876$ & $0.871$ & $0.707$ & $+1.13$ \\
    Energy  & $0.905$ & $0.859$ & $0.618$ & $+1.40$ \\
    Tension & $0.945$ & $0.913$ & $0.736$ & $+1.23$ \\
    Anger   & $0.922$ & $0.919$ & $0.918$ & $+0.10$ \\
    Fear    & $0.933$ & $0.931$ & $0.929$ & $+0.17$ \\
    Happy   & $0.930$ & $0.926$ & $0.926$ & $+0.06$ \\
    Sad     & $0.877$ & $0.869$ & $0.822$ & $+0.60$ \\
    Tender  & $0.908$ & $0.891$ & $0.886$ & $+0.21$ \\
    \hline
    Mean    & $0.912$ & $0.897$ & $0.818$ & $+0.61$ \\
    \hline
  \end{tabular}
\end{table}

Agreement is high throughout (mean $r=0.912$), but the two intraclass coefficients
diverge: consistency averages $0.897$ while absolute agreement averages only $0.818$, and
the difference is confined to the three bipolar scales -- valence, energy and tension --
where the second panel rated systematically about $1.25$ points higher
($d=0.68$--$0.89$, $p<10^{-23}$). On the five discrete emotion intensities the two panels
differ by $0.06$--$0.21$ points, mostly not significantly. The panels therefore agree
closely about \emph{what} they hear and differ in \emph{how they use the response scale}.
This distinction matters for the modelling: after standardising each set independently --
which is what the classifier pipeline does to its inputs -- the mean absolute difference
falls to $0.305$ standard deviations, so the shift is absorbed by the preprocessing rather
than propagating into the results.

Excerpts presented twice to the \emph{same} panel are reproduced almost exactly (mean
$r=0.988$ over the seven repeat trials, with negligible bias), indicating that the
excerpts themselves are not intrinsically ambiguous and that the between-panel difference
reflects panel and context rather than uncertainty about the stimuli. This estimate rests
on only seven excerpts and on repetitions within a single session, so it should be read as
an upper bound on within-panel reliability.

\paragraph{Consequences for the emotion regression.}
A predictor cannot correlate with a noisy target more strongly than the target correlates
with itself. Treating the between-panel consistency as a parallel-forms reliability
$r_{xx}$, the largest attainable coefficient of determination is $r_{xx}$ itself, here
$0.897$ on average. Consistency rather than absolute agreement is the appropriate
quantity, since the regression is scored against a single set's ratings and a constant
offset between panels does not reduce the variance explainable within that set. The random
forest reaches a mean $R^2$ of $0.560$ (Section~\ref{subsec:emotion_regression}), that is
about $62\,\%$ of the attainable variance.

Two conclusions follow. First, a coefficient of determination near unity was never
achievable, so $0.560$ should be read against a ceiling of roughly $0.90$ rather than
against $1.0$. Second, and less comfortably, rating noise does \emph{not} account for the
shortfall: with more than a third of the attainable variance still unexplained, the
limiting factor is the audio representation and the amount of training data rather than
the quality of the labels. This agrees with the learning-curve analysis in
Section~\ref{subsec:learning_curve}, where performance had not plateaued at the full
corpus size, and it is consistent with the range reported for music emotion recognition
generally \citep{kang2025there}.
